\chapter{Data and Graph Construction Methodology}
\label{chap:method}

\section{Corpus definition and data sources}

The release corpus is the frozen roster-v10 set of AAAI and IJCAI main-track
papers published from 2021 to 2025. The official proceedings pages establish
whether a paper is in scope, the edition year, title, byline order, and official
record URL. The roster contains \num{13617} papers: \num{9910} AAAI and \num{3707}
IJCAI papers. This two-series release is a subset of a wider source registry and
must not be interpreted as a census of AI research. The official source portals
are listed in \cref{chap:reproduction} and the bibliography
\citep{aaaiProceedings,ijcaiProceedings}.

OpenAlex provides matched work metadata, author records, institution links,
abstracts, and references. The API is accessed through a source cache that stores
response bytes, request metadata, retrieval times, and SHA-256 hashes
\citep{openalexapi,openalexworks}. Raw downloads are treated as immutable inputs;
the report package does not redistribute the full cache or papers. The source
roles remain distinct: official rosters govern inclusion and printed bylines,
while OpenAlex is an enrichment source subject to matching and identity checks.

\section{Acquisition, matching, and duplicate handling}

For each roster record, the acquisition code derives a DOI from the official
record URL using venue-specific rules. It requests OpenAlex works by DOI rather
than searching by title. A candidate work is accepted only if its returned DOI
agrees exactly, its normalized title similarity is at least 0.90, and its
authorship count matches the official byline count. Title similarity uses
normalized strings and a sequence-matching ratio; a near match is retained as a
flag for inspection. When multiple works share a DOI, the official landing-page
URL can resolve the duplicate only if exactly one matching candidate passes the
checks. Otherwise the paper remains in review.

The measured acquisition matched \num{13606} of \num{13617} roster papers
(99.92\%). Eleven records remain in review: nine duplicate-DOI cases without an
unambiguous eligible work and two author-count mismatches. No paper was reported
as not found after the acquisition retries. All matched works contain an abstract;
two accepted matches carry a title-similarity flag. These counts describe work
matching, not author identity accuracy. Re-running a current API query is not
guaranteed to reproduce the original bytes, so the versioned cache hashes are
part of the reproducibility record.

\section{Author and researcher resolution}

Each official author position is compared with the corresponding OpenAlex
authorship in the DOI-matched work. A researcher node is keyed by the OpenAlex
author ID only when the position is compatible with the official name under the
release rule. ORCID is stored when present but is not required; names alone never
create merges. Positions that lack an author ID, have an incompatible name, or
participate in a paper-level conflict are withheld rather than forced into a
canonical person record. Stable internal IDs and external-ID crosswalks preserve
the distinction between graph identity and display names.

The final correction that produced v4 checks both singular ORCID values and
retained observed-ORCID arrays across every matched work. Multiple distinct,
checksum-valid ORCIDs for one provider author ID trigger global withholding;
missing or null ORCID remains an absence of evidence. The scan found \num{4492}
conflicting OpenAlex author IDs and no malformed nonempty ORCID values. It removed
\num{8722} previously accepted authorships, \num{4491} researcher nodes, and 49
organizations with no surviving affiliation observations. One conflicting
provider ID was already absent from the graph. The method is a versioned
correction, not a change to official papers, topic labels, or citation claims.

Of \num{65507} official author positions, \num{65469} occur on matched works.
The corrected graph contains \num{50437} accepted AUTHORED relationships and
withholds \num{15032} matched positions. There are 188 matched papers with no
accepted authorship and 199 across the entire roster, including the 11 unmatched
papers. These missing positions constrain the researcher projection and can
change components and centrality rankings.

\section{Affiliation evidence and organization resolution}

Affiliations are first retained as paper-specific authorship observations. Each
observation ties a particular author position on a particular paper to an
OpenAlex institution ID and its raw affiliation string. Only after these triples
are preserved are the graph's researcher--organization edges aggregated. This
prevents the paper-level institution list from being copied to every author and
avoids treating a current profile as evidence of historical affiliation.

The evidence supports \num{48259} institution assignments among \num{50437}
accepted authorships. That is 95.68\% conditional coverage of accepted authors,
but the course coverage denominator includes withheld matched byline positions:
\num{48259}/\num{65469}=73.71\%. Against all official positions the ratio is
73.67\%. Both project-level measures fall below the 80\% target. The conditional
percentage is useful as a source diagnostic but does not replace the coverage
gate. Organization sector and country are normalized organization metadata;
unknown values remain explicitly unknown.

\section{Topic classification}

The project taxonomy has four parent topics and 16 leaf topics. The frozen
classifier encodes topic definitions and paper title/abstract text using
\modelid{sentence-transformers/all-MiniLM-L6-v2}, pinned at revision
\texttt{1110a243fdf4706b3f48f1d95db1a4f5529b4d41}. It applies the development
calibration thresholds $\tau=0.2708871365$ and $\delta=0.02$, assigns at most
three leaf topics, and abstains when its rule does not support a label. A source
sentence from the abstract, or the title for title-only input, is retained as
topic evidence. The model and its frozen revision are documented with the
project's topic configuration and model card \citep{reimers2019sentencebert,minilmModelCard}.

The corrected graph retains \num{11575} topic-labeled papers and \num{2042}
abstentions, for 85.00\% coverage. It contains \num{14222} \texttt{HAS\_TOPIC} edges
because some papers have multiple leaf labels. The label count is not an
accuracy score. Thresholds were calibrated on development papers; the reserved
test IDs were not used to tune the classifier. The only reported topic test
metrics use historical AI-assigned labels and are discussed separately in
\cref{chap:evaluation}.

\section{Graph model and edge semantics}

The graph has five entity labels and ten predicates. Observed bibliographic
relationships are kept distinct from model-inferred topic links, curated
taxonomy relationships, and social/topic edges derived from accepted core-paper
evidence. The schema table in \cref{tab:relationship-schema} gives each
predicate's direction, evidence class, and support rule. In aggregate, every
edge has a deterministic identity based on predicate and endpoints, accepted
status, evidence IDs, method version, and graph version.

\begin{landscape}
\begin{table}[p]
  \centering
  \caption{The ten graph predicates, their direction, evidence class, and release
  support rule. Symmetric collaboration predicates are stored once in canonical
  order and analyzed as undirected.}
  \label{tab:relationship-schema}
  \small
  \begin{tabularx}{\linewidth}{>{\ttfamily\raggedright\arraybackslash}p{3.55cm}p{4.3cm}p{2.0cm}X}
    \toprule
    \textbf{Predicate} & \textbf{Source $\rightarrow$ target} &
    \textbf{Class} & \textbf{Meaning and acceptance rule}\\
    \midrule
    AUTHORED & Researcher $\rightarrow$ Paper & Observed & Compatible resolved identity at the official byline position; preserve position and roster name.\\
    AFFILIATED\_WITH & Researcher $\rightarrow$ Organization & Observed & Accepted author--paper--organization observations, aggregated only after paper-specific triples are preserved; not employment duration.\\
    PUBLISHED\_IN & Paper $\rightarrow$ Venue & Observed & Official in-scope proceedings record, one primary series per core paper.\\
    CITES & Paper $\rightarrow$ Paper & Observed & Explicit reference to a verified, deduplicated target; direction is citing to cited.\\
    HAS\_TOPIC & Paper $\rightarrow$ Topic & Inferred & One to three supported leaf labels from the frozen classifier, or abstention.\\
    SUBTOPIC\_OF & Topic $\rightarrow$ Topic & Curated & Frozen approved leaf-to-parent taxonomy relation.\\
    COAUTHORED\_WITH & Researcher $\leftrightarrow$ Researcher & Derived & Distinct resolved people share an accepted core paper; fractional contribution uses the full byline denominator.\\
    COLLABORATED\_WITH & Organization $\leftrightarrow$ Organization & Derived & One core paper supports two organizations through two distinct authors with aligned paper-specific affiliations; a sole dual-affiliated author is insufficient.\\
    RESEARCHES\_TOPIC & Researcher $\rightarrow$ Topic & Derived & At least two distinct accepted core papers by the researcher carry the leaf topic.\\
    \makecell[l]{\texttt{CONTRIBUTES\_TO\_}\\\texttt{TOPIC}} & Organization $\rightarrow$ Topic & Derived & At least three distinct accepted core papers carry the leaf topic and have paper-specific organization support.\\
    \bottomrule
  \end{tabularx}
\end{table}
\end{landscape}


\begin{figure}[tbp]
\centering
\begin{tikzpicture}[
  font=\small,
  source/.style={draw=UETBlue,fill=PaleBlue,rounded corners=3pt,align=center,text width=5.0cm,minimum height=1.05cm,inner sep=6pt},
  stage/.style={draw=Slate,fill=white,rounded corners=3pt,align=center,text width=9.0cm,minimum height=1.1cm,inner sep=7pt},
  arrow/.style={-{Latex[length=2.2mm]},thick,color=UETBlue},
  node distance=7mm and 10mm
]
  \node[source] (aaai) {Official AAAI proceedings roster\\2021--2025 main track};
  \node[source,right=of aaai] (ijcai) {Official IJCAI proceedings roster\\2021--2025 main track};
  \node[source,below=of ijcai] (openalex) {OpenAlex DOI work records\\cached response bytes and source hashes};
  \node[stage,below=14mm of $(aaai)!0.5!(ijcai)$] (match) {DOI-derived work matching\\exact DOI, title similarity $\geq 0.90$, and author-count agreement};
  \node[stage,below=of match] (identity) {Position-specific author resolution and affiliation triples\\compatible names; stable author IDs; conflicting ORCID evidence withheld};
  \node[stage,below=of identity] (topic) {Frozen text classification and taxonomy\\title/abstract embeddings; one to three leaf labels or abstention};
  \node[stage,below=of topic] (graph) {Typed graph and evidence ledger\\observed R1--R4; inferred R5; curated R6; derived R7--R10};
  \node[stage,below=of graph] (validate) {Versioned validation and analysis\\structural checks, size gates, named projections, figures, and source-linked findings};
  \draw[arrow] (aaai.south) -- (match.north west);
  \draw[arrow] (ijcai.south) -- (match.north);
  \draw[arrow] (openalex.west) -| (match.east);
  \draw[arrow] (match) -- (identity);
  \draw[arrow] (identity) -- (topic);
  \draw[arrow] (topic) -- (graph);
  \draw[arrow] (graph) -- (validate);
\end{tikzpicture}
\caption{Data-to-graph workflow for the frozen release. Metadata matching,
identity resolution, text classification, and relationship derivation are separate
steps. The official roster defines inclusion; OpenAlex enriches matched works.}
\label{fig:data-pipeline}
\end{figure}


\section{Derived relationships and weights}

Derived relationships are materialized only from accepted base evidence. For a
core paper with $k>1$ byline authors, each eligible resolved researcher pair
receives fractional coauthorship weight $1/(k-1)$; unresolved authors remain in
the denominator so that missing identities do not inflate accepted pair weights.
The distinct shared-paper count is retained alongside the fractional strength.

For organization collaboration, the construction enumerates distinct
organization pairs supported by distinct authors with paper-specific
affiliations. A sole author with two affiliations creates no collaboration edge.
If a paper supports $q>0$ eligible organization pairs, each receives weight
$1/q$ for that paper; support is deduplicated by paper. Researcher--topic edges
require at least two distinct accepted core papers, and organization--topic edges
require at least three distinct papers with the topic and a paper-specific
organization support. Topic and organization portfolios divide each labeled
paper's mass fractionally across its accepted labels and, where appropriate,
supported organizations. Those portfolio weights do not replace the raw-paper
thresholds used to create R9 and R10.

\section{Network projections}

The analysis keeps five projections separate: undirected researcher
coauthorship, undirected organization collaboration, directed in-corpus paper
citations, paper--topic membership, and organization--topic contributions.
Collaboration pairs are stored once in canonical ID order and analyzed as
undirected relationships. Citation arrows point from citing paper to cited
paper. Weighted shortest-path calculations use the inverse of positive
fractional edge strength. The results report projection boundaries and seeds so
that a number such as betweenness is interpretable only under its stated graph
and distance rule.
